\documentclass[12pt]{article}

\usepackage{amsmath, amsthm, amssymb}
\usepackage{geometry}
\usepackage[colorlinks=true, linkcolor=blue, citecolor=blue, urlcolor=blue]{hyperref}
\usepackage{microtype}

\geometry{margin=1.25in}

\newtheorem{theorem}{Theorem}
\newtheorem{lemma}[theorem]{Lemma}
\newtheorem{proposition}[theorem]{Proposition}
\newtheorem{corollary}[theorem]{Corollary}
\theoremstyle{definition}
\newtheorem{definition}[theorem]{Definition}
\newtheorem{property}[theorem]{Property}

\title{Suggested Revisions to Chapter~24:\\
\textit{Credible Government Policies: I}\\[0.5ex]
\large \textit{Recursive Macroeconomic Theory}, 4th edition}
\author{Based on the note by Barth\'{e}lemy and Mengus (2020)}
\date{}

\begin{document}

\maketitle

\begin{abstract}
\noindent
This document collects suggested revisions to the statements and proofs in
Chapter~24 of Ljungqvist and Sargent (2018). The core issue, identified by
Barth\'{e}lemy and Mengus (2020), is that the two-value characterization of
supported equilibria stated on page~1034 is presented as a necessary and
sufficient condition, but is in fact only sufficient at the point where it is
invoked. It becomes necessary only \emph{after} $V$ is known to be compact,
which is precisely what the surrounding argument is trying to establish. The
fix requires replacing every admissible \emph{4-tuple} $(x, y, w_1, w_2)$ with
an admissible \emph{triple} $(x, y, w(\cdot))$ where $w(\cdot)$ is a function
on $Y$, proving compactness of $V$ using this refined machinery, and then
recovering the two-value formulation as a corollary. The economic content of
the chapter is entirely unaffected.
\end{abstract}

\tableofcontents

\bigskip

%------------------------------------------------------------------
\section{The Gap and Its Source}
%------------------------------------------------------------------

On page~1034, the text states that a competitive equilibrium $(x, y) \in C$
is supported by a subgame perfect equilibrium (SPE)---that is, $x = \sigma^h_1$
and $y = \sigma^g_1$---\emph{if and only if} there exist two values
$(v_1, v_2) \in V \times V$ such that
\begin{equation}
(1-\delta)\,r(x,y) + \delta v_1
\;\geq\;
(1-\delta)\,r(x,\eta) + \delta v_2,
\qquad \forall\, \eta \in Y.
\label{eq:two-value}
\end{equation}
No explicit proof of the ``only if'' direction is provided. As we show below,
condition~\eqref{eq:two-value} is \emph{sufficient} but not \emph{necessary}
unless $V$ is already known to be compact---but compactness of $V$ is precisely
what the subsequent argument aims to establish.

The correct necessary and sufficient condition replaces the two scalar
continuation values $(v_1, v_2)$ with a \emph{function} $v(\cdot) : Y \to V$,
one continuation value for each possible government action. We state and prove
this refined lemma in Section~\ref{sec:lemma}, use it to rebuild the APS
machinery in Sections~\ref{sec:def7}--\ref{sec:theorem1}, prove compactness of
$V$ in Section~\ref{sec:compact}, and then recover the two-value
characterization as a corollary in Section~\ref{sec:corollary}.

%------------------------------------------------------------------
\section{Revised Characterization of Supported Equilibria}
\label{sec:lemma}
%------------------------------------------------------------------

The following lemma replaces the ``if and only if'' assertion on page~1034.

\begin{lemma}[Necessary and sufficient condition for support by an SPE]
\label{lem:nsc}
A competitive equilibrium $(x, y) \in C$ is supported by an SPE, i.e.\
$x = \sigma^h_1$ and $y = \sigma^g_1$, if and only if there exists a function
$v(\cdot) : Y \to V$ such that
\begin{equation}
(1-\delta)\,r(x,y) + \delta v(y)
\;\geq\;
(1-\delta)\,r(x,\eta) + \delta v(\eta),
\qquad \forall\, \eta \in Y.
\label{eq:func-cond}
\end{equation}
\end{lemma}

\begin{proof}
\textbf{Necessity.}
Suppose $(x, y) \in C$ is supported by an SPE $\sigma$, so that
$x = \sigma^h_1$ and $y = \sigma^g_1$.
By condition~(1) of the Lemma on page~1030, for each $\eta \in Y$ the
continuation strategy $\sigma|_{(x,\eta)}$ is itself an SPE.
Define $v(\eta) = V_g(\sigma|_{(x,\eta)}) \in V$ for each $\eta \in Y$.
Condition~(3) of the Lemma on page~1030 then gives
\[
(1-\delta)\,r(x,y) + \delta V_g(\sigma|_{(x,y)})
\;\geq\;
(1-\delta)\,r(x,\eta) + \delta V_g(\sigma|_{(x,\eta)}),
\]
which is exactly \eqref{eq:func-cond} with $v(\eta) = V_g(\sigma|_{(x,\eta)})$.

\textbf{Sufficiency.}
Suppose there exists $v(\cdot) : Y \to V$ satisfying \eqref{eq:func-cond}.
Since $v(\eta) \in V$ for each $\eta \in Y$, there exists for each $\eta$ an
SPE $\tilde\sigma(\eta)$ with $V_g(\tilde\sigma(\eta)) = v(\eta)$.
Construct a strategy profile $\sigma$ by setting
$\sigma|_{(x,\eta)} = \tilde\sigma(\eta)$ for all $\eta \in Y$, and
(without loss of generality) $\sigma|_{(\nu,\eta)} = \tilde\sigma(\eta)$
for all $\nu \in X$.
By construction, this profile satisfies all three conditions of the Lemma on
page~1030 and is therefore an SPE with $\sigma^h_1 = x$ and $\sigma^g_1 = y$.
\end{proof}

\bigskip

The two-value condition~\eqref{eq:two-value} of the original text is
\emph{sufficient}: if $(v_1, v_2) \in V \times V$ satisfy it, set $v(y) = v_1$
and $v(\eta) = v_2$ for $\eta \neq y$, and \eqref{eq:func-cond} holds.
The reverse implication holds once $V$ is known to be compact
(see Section~\ref{sec:corollary}), but cannot be assumed at this stage.

%------------------------------------------------------------------
\section{Revised Definition 7: Admissible Triples}
\label{sec:def7}
%------------------------------------------------------------------

Using Lemma~\ref{lem:nsc}, Definition~7 on page~1034 should be replaced as
follows.

\begin{definition}[Admissible triple]
\label{def:admissible}
Let $W \subset \mathbb{R}$.  A triple $(x, y, w(\cdot))$ is said to be
\emph{admissible with respect to $W$} if $(x, y) \in C$ and
$w(\cdot) : Y \to W$ is a function satisfying
\[
(1-\delta)\,r(x,y) + \delta w(y)
\;\geq\;
(1-\delta)\,r(x,\eta) + \delta w(\eta),
\qquad \forall\, \eta \in Y.
\]
\end{definition}

%------------------------------------------------------------------
\section{Revised Definition 8: The Operator $B$}
\label{sec:def8}
%------------------------------------------------------------------

Definition~8 on page~1036 is updated to reflect the triple-based notion of
admissibility.

\begin{definition}[The operator $B$]
\label{def:B}
For each set $W \subset \mathbb{R}$, let $B(W)$ be the set of possible values
\[
w \;=\; (1-\delta)\,r(x,y) + \delta\,w(y)
\]
associated with admissible triples $(x, y, w(\cdot))$ with respect to $W$.
\end{definition}

%------------------------------------------------------------------
\section{Revised Monotonicity of $B$}
\label{sec:mono}
%------------------------------------------------------------------

The monotonicity property is unchanged in statement; only the proof is
adjusted.

\begin{property}[Monotonicity of $B$]
If $W \subseteq W' \subseteq \mathbb{R}$, then $B(W) \subseteq B(W')$.
\end{property}

\begin{proof}
Suppose $w \in B(W)$.  Then there exists an admissible triple $(x, y, w(\cdot))$
with respect to $W$ such that $w = (1-\delta)r(x,y)+\delta w(y)$.
Since $W \subseteq W'$, the function $w(\cdot)$ maps $Y$ into $W'$ as well, so
$(x, y, w(\cdot))$ is admissible with respect to $W'$ and $w \in B(W')$.
\end{proof}

%------------------------------------------------------------------
\section{Revised Theorem 1: Self-Generating Sets Are Subsets of $V$}
\label{sec:theorem1}
%------------------------------------------------------------------

Definition~9 (self-generating sets) is unchanged: $W$ is self-generating if
$W \subseteq B(W)$.  Theorem~1 is restated and its proof is adapted to use
admissible triples.

\begin{theorem}[Self-generating sets are subsets of $V$]
\label{thm:selfgen}
If $W \subset \mathbb{R}$ is bounded and self-generating, then $B(W) \subseteq V$.
\end{theorem}

\begin{proof}
Assume $W \subseteq B(W)$.  Choose $w \in B(W)$ and construct an SPE with
value $w$ as follows.

\smallskip
\noindent\textbf{Step 1.}
Because $w \in B(W)$, there exist $(x_1, y_1) \in C$ and a function
$w_1(\cdot) : Y \to W$ such that
\[
w
\;=\;
(1-\delta)\,r(x_1, y_1) + \delta\,w_1(y_1)
\;\geq\;
(1-\delta)\,r(x_1,\eta) + \delta\,w_1(\eta),
\qquad \forall\, \eta \in Y.
\]
Set $\sigma_1 = (x_1, y_1)$.

\smallskip
\noindent\textbf{Step 2.}
Since $w_1(y_1) \in W \subseteq B(W)$, there exist $(x_2, y_2) \in C$ and a
function $w_2(\cdot) : Y \to W$ such that
\[
w_1(y_1)
\;=\;
(1-\delta)\,r(x_2, y_2) + \delta\,w_2(y_2)
\;\geq\;
(1-\delta)\,r(x_2,\eta) + \delta\,w_2(\eta),
\qquad \forall\, \eta \in Y.
\]
Set the first-period outcome of the continuation economy (occurring given that
$y_1$ was chosen in period~1) equal to $(x_2, y_2)$; that is, set
$(\sigma|_{(x_1,y_1)})_1 = (x_2, y_2)$.

\smallskip
\noindent\textbf{Continuation.}
Iterating in this way, for each $w \in B(W)$ we construct sequences of
functions $w_1(\cdot), w_2(\cdot), w_3(\cdot), \ldots$ and first-period
outcomes $(x_1, y_1), (x_2, y_2), (x_3, y_3), \ldots$.
Since $\delta < 1$ and $W$ is bounded,
\[
w
\;=\;
\lim_{t\to\infty}
\left[
  (1-\delta)\sum_{s=1}^{t} \delta^{s-1} r(x_s, y_s)
  + \delta^t\, w_t(y_t)
\right]
\;=\;
(1-\delta)\sum_{s=1}^{\infty} \delta^{s-1} r(x_s, y_s).
\]
At each stage, the continuation values make one-period deviations
unimprovable, so the constructed strategy profile is an SPE.  By construction,
$V_g(\sigma) = w$, which shows $w \in V$.
\end{proof}

%------------------------------------------------------------------
\section{Compactness of $V$}
\label{sec:compact}
%------------------------------------------------------------------

With the revised machinery in place, compactness of $V$ follows by an argument
parallel to the proof of Theorem~4 in Abreu, Pearce, and Stacchetti (1990).

\begin{proposition}[Compactness of $V$]
\label{prop:compact}
The set $V$ is compact.
\end{proposition}

\begin{proof}
We know that $V$ is bounded and self-generating ($V \subseteq B(V)$).
Let $\mathrm{cl}(V)$ denote the closure of $V$.  Since $V$ is bounded,
$\mathrm{cl}(V)$ is compact and $V \subseteq \mathrm{cl}(V)$.

By monotonicity of $B$,
\[
V \;=\; B(V) \;\subseteq\; B(\mathrm{cl}(V)).
\]
Since $B(\cdot)$ maps compact sets into compact sets, $B(\mathrm{cl}(V))$ is
compact.  Because $\mathrm{cl}(V)$ is the smallest closed (hence compact) set
containing $V$, and $V \subseteq B(\mathrm{cl}(V))$, it follows that
\[
\mathrm{cl}(V) \;\subseteq\; B(\mathrm{cl}(V)),
\]
so $\mathrm{cl}(V)$ is bounded and self-generating.
Theorem~\ref{thm:selfgen} then gives $\mathrm{cl}(V) \subseteq V$, and
therefore $V = \mathrm{cl}(V)$, which means $V$ is compact.
\end{proof}

%------------------------------------------------------------------
\section{Corollary: The Two-Value Characterization}
\label{sec:corollary}
%------------------------------------------------------------------

Now that $V$ is known to be compact, the function-valued condition of
Lemma~\ref{lem:nsc} collapses to the economically transparent two-value form,
and the assertion on page~1034 can be stated cleanly as the following
corollary.

\begin{corollary}[Two-value characterization]
\label{cor:two-value}
A competitive equilibrium $(x, y) \in C$ is supported by an SPE, that is
$x = \sigma^h_1$ and $y = \sigma^g_1$, if and only if there exist
$(v_1, v_2) \in V \times V$ such that
\[
(1-\delta)\,r(x,y) + \delta v_1
\;\geq\;
(1-\delta)\,r(x,\eta) + \delta v_2,
\qquad \forall\, \eta \in Y.
\]
\end{corollary}

\begin{proof}
\textbf{Sufficiency} is immediate: given $(v_1, v_2) \in V \times V$, define
$v(y) = v_1$ and $v(\eta) = v_2$ for $\eta \neq y$; then
\eqref{eq:func-cond} holds, so Lemma~\ref{lem:nsc} applies.

\textbf{Necessity.}
Since $(x,y)$ is supported by an SPE, Lemma~\ref{lem:nsc} yields a function
$v(\cdot) : Y \to V$ satisfying \eqref{eq:func-cond}.
Set $v_1 = v(y)$.
The set $\{v(\eta) : \eta \in Y\}$ is a subset of $V$; since $V$ is compact
(Proposition~\ref{prop:compact}) and bounded, its infimum is attained in $V$.
Set
\[
v_2 \;=\; \inf_{\eta \in Y}\, v(\eta) \;\in\; V.
\]
Then $v(\eta) \geq v_2$ for all $\eta \in Y$, so \eqref{eq:func-cond} implies
\[
(1-\delta)\,r(x,y) + \delta v_1
\;=\;
(1-\delta)\,r(x,y) + \delta v(y)
\;\geq\;
(1-\delta)\,r(x,\eta) + \delta v(\eta)
\;\geq\;
(1-\delta)\,r(x,\eta) + \delta v_2,
\]
for all $\eta \in Y$, as required.
\end{proof}

%------------------------------------------------------------------
\section{Summary of Changes}
%------------------------------------------------------------------

The table below lists the specific replacements needed in the text of
Chapter~24.

\medskip
\begin{center}
\begin{tabular}{lll}
\hline
\textbf{Location} & \textbf{Current text} & \textbf{Replacement} \\
\hline
p.\ 1034, main assertion & ``if and only if $\exists\,(v_1,v_2)$\ldots'' &
  Lemma~\ref{lem:nsc} above \\
p.\ 1034, Definition~7 & admissible 4-tuple $(x,y,w_1,w_2)$ &
  Definition~\ref{def:admissible} above \\
p.\ 1036, Definition~8 & $B(W)$ via 4-tuples &
  Definition~\ref{def:B} above \\
p.\ 1036, Monotonicity proof & uses $(w_1,w_2)$ pair &
  Section~\ref{sec:mono} above \\
p.\ 1037, Theorem~1 proof & uses scalar $w_1, w_2$ &
  Theorem~\ref{thm:selfgen} above \\
p.\ 1038, after Theorem~1 & compactness asserted without proof &
  Proposition~\ref{prop:compact} above \\
p.\ 1034, ``if and only if'' & stated without full proof &
  Corollary~\ref{cor:two-value} above \\
\hline
\end{tabular}
\end{center}

\medskip
\noindent
All other material in Chapter~24---the economic motivation, the examples
(infinite repetition of Nash, trigger strategies, when reversion to Nash is not
bad enough), the iteration algorithm for computing $V$, and the discussion of
best and worst SPE values---is unaffected by these changes.

%------------------------------------------------------------------
\begin{thebibliography}{9}
\bibitem{APS1990}
Abreu, D., D.\ Pearce, and E.\ Stacchetti (1990):
``Toward a Theory of Discounted Repeated Games with Imperfect Monitoring,''
\textit{Econometrica}, 58, 1041--1063.

\bibitem{BM2020}
Barth\'{e}lemy, J.\ and E.\ Mengus (2020):
``Note on Chapter~24 on Credible Government Policies: Proof of the Compactness
of the Set of Equilibrium Payoffs,'' unpublished manuscript.

\bibitem{LS2018}
Ljungqvist, L.\ and T.\ J.\ Sargent (2018):
\textit{Recursive Macroeconomic Theory}, 4th edition, MIT Press.
\end{thebibliography}

\end{document}
